% !TEX root = ../manuscript.tex
\section{Knowledge Engineering Representation Checks}
\label{sec:application}
\label{sec:experimental-results}

This section evaluates fixed-budget audit as an information-representation task. The experiments ask whether audit records track annotation signals, whether graph construction changes dependency representations, and whether decomposed records expose audit-target information under budget constraints.

\subsection{Evaluation Philosophy: Representation Checks}
\label{sec:evaluation-philosophy}

The metrics in this section are representation checks. Spearman $\rho$ and Kendall $\tau$ measure whether an audit-priority field tracks the ordering of a supplied annotation or fidelity field. NDCG, Coverage@25\%, Recall, first-selected target, and Mass@25\% measure whether relevant audit-record information remains visible under a fixed review budget. These readouts evaluate the structure and visibility of the record.

The reported values indicate fidelity tracking, coverage, or visibility for the specified representation object under the specified data source.

\subsection{Task Definition: Fixed-Budget Knowledge Audit}
\label{sec:application-setup}

\subsubsection{Data Sources and Reference Views}

The process-annotation stage uses PRM800K phase2 records under the frozen v3.6/v3.8 hash split. A deterministic SHA-256 split over a 12,000-row source pool assigns 4,344 samples (33,555 labeled artifacts) to development and 4,417 samples (34,219 labeled artifacts) to locked reporting. The remaining 3,239 rows are retained in the split manifest but fail eligibility filters for labeled artifact coverage, split assignment, or archived feature completeness. The \wstruct{} Ridge model is fit only on development and frozen before locked reporting; the regression target is a PRM800K completion rating selected by the deterministic extraction rule: use the recorded chosen completion when it has an unflagged rating, otherwise use the first unflagged rated completion. This rating is linearly mapped to $[0,1]$ as $(\mathrm{rating}+1)/2$ and clipped to the unit interval. Its feature policy forbids label fields and raw completion ratings, but the model remains supervised and uses engineered lexical, positional, numeric, and cue features. Comparisons between \wstruct{} and unsupervised structural fields are therefore not comparisons under equal supervision; the current package does not include a same-supervision structure-only Ridge control. The Countries-KG stage tests backend feasibility. The controlled synthetic stage uses 200 traces over five fixed seeds for calibration, ablation, and error analysis. MuSiQue, WebQSP, and failed GSM8K/HotpotQA routes remain supplementary boundary checks.

The frozen \wstruct{} feature list is recorded here to make the fidelity-field construction reproducible and mechanically checkable against the archived model:
% w_struct_feature_list_begin
\begin{center}
\scriptsize
\begin{tabular}{lll}
\texttt{raw\_local\_utility} & \texttt{relative\_position} & \texttt{relative\_position\_sq} \\
\texttt{relative\_position\_cu} & \texttt{log\_token\_count} & \texttt{numeric\_density} \\
\texttt{equation\_density} & \texttt{answer\_cue} & \texttt{conclusion\_cue\_count} \\
\texttt{error\_uncertainty\_cue\_count} & \texttt{reasoning\_cue\_count} & \texttt{is\_first\_step} \\
\texttt{is\_last\_step} & \texttt{trace\_step\_count} & \texttt{source\_step\_index\_ratio}
\end{tabular}
\end{center}
% w_struct_feature_list_end

Reference views include raw local utility, span length, relative position, seeded random field assignment, and stage-specific scalar views. The primary readout is annotation-order tracking for the audit-priority field, measured by Spearman $\rho$; Kendall $\tau$, record-level NDCG, and fixed-budget coverage are reported where they clarify representation visibility. Sensitivity grids, additional ablations, error modes, token-level attribution scope, and efficiency measurements are in the supplementary material.

The controlled synthetic stage tests calibration under skewed local utility, sparse necessity, and block-structured redundancy. Its annotation target combines correctness, lexical diversity, position, and structural information, which makes the audit field depend on more than one SCU input vector.

For the larger real-data stage, PRM800K is treated as a process-annotation distribution. The locked v3.6 split examines annotation-order consistency under a fixed audit budget. The v3.8 frozen PRM reference provides in-distribution prefix-sequence context. The resulting evidence concerns process-annotation audit-record behavior; deployment-oriented RAG, KG-maintenance, and human-curation workflows are discussed as future validation settings.

All reported readouts are deterministic under archived configs. The controlled synthetic stage uses five fixed seeds; confidence intervals, extended readouts, tests, and effect sizes are retained in the supplementary material.

\subsubsection{Evidence Stages and Interpretation}

The evidence is staged because each dataset tests a different part of the workflow. Table~\ref{tab:evidence-routes} states the interpretation attached to each stage.

\begin{table}[ht]
\caption{Evidence stages and interpretation. Each stage instantiates a separate part of the audit-record construction process.}
\label{tab:evidence-routes}
\centering
\small
\renewcommand{\arraystretch}{1.12}
\begin{tabular}{Z{0.27\linewidth}Z{0.27\linewidth}Z{0.34\linewidth}}
\toprule
\textbf{Evidence stage} & \textbf{Primary evidence object} & \textbf{Manuscript interpretation} \\
\midrule
PRM800K process-annotation stage & 4,417 samples, 34,219 labeled artifacts & Process-annotation consistency and audit context. \\
\makecell[l]{Countries-KG backend\\feasibility study} & \makecell[l]{Fixture-level\\Countries-KG\\check} & \makecell[l]{Knowledge-graph backend\\feasibility and typed\\structural-dependency sensitivity.} \\
Audit-target retrieval stage & Rule-derived audit targets from process annotations & Lowest-overlap audit-target visibility, with aggregate context exploratory. \\
Controlled synthetic stage & 200 traces, approximately 1,000 steps, five fixed seeds & Controlled calibration and ablation evidence. \\
Frozen PRM reference & Prefix-sequence reward signals & In-distribution reference context. \\
Supplementary MuSiQue/WebQSP checks & Supplementary details only & Feasibility and fixed-schema context only. \\
Failed GSM8K/HotpotQA attempts & Failed or blocked runs & Transparency only; no positive evidence claim. \\
\bottomrule
\end{tabular}
\end{table}

This staged interpretation also governs null or negative results. Full QP can be too aggressive when the base fidelity field already aligns with annotation structure. The contribution is reproducible representation under stated data-generating conditions.

\subsection{Representation Fidelity Tracking on Process Annotations}
\label{sec:prm800k-evidence}

The main real-data fidelity check uses the locked PRM800K side for annotation-field tracking and fixed-budget audit-record coverage. The stored \wstruct{} field is the primary fidelity field (Spearman $\rho = 0.611$), computed by a dev-trained frozen Ridge model. Raw local utility is negative ($\rho = -0.077$), and the overlap-limited frozen PRM prefix field is lower ($\rho = 0.252$). SC-FMA Ridge closely tracks \wstruct{} ($\rho = 0.604$), while QP ($\rho = 0.442$) and Projection ($\rho = -0.135$) are lower-consistency variants. The archived paired-difference report does not support treating any SC-FMA variant as an annotation-order improvement over \wstruct{} on this route: Ridge differs by $-0.007$ Spearman (bootstrap CI $[-0.011,-0.004]$), QP by $-0.169$ (CI $[-0.176,-0.163]$), and Projection by $-0.746$ (CI $[-0.760,-0.732]$). This is not evidence that the redundancy and bottleneck terms improve PRM800K ranking; those terms are interpreted as structural checks and conditional regularizers. Projection is the clearest failure mode: the topology-constrained simplex step can reorder artifacts when structural constraints conflict with a strong fidelity field, so it is kept as a fallback check; Ridge remains the conservative reporting variant. Supplementary Tables C.8 and C.9 report confidence intervals and fixed-budget readouts.

\begin{table}[ht]
\caption{PRM800K annotation-order consistency and fixed-budget audit coverage on the frozen split of 4,417 samples and 34,219 labeled steps. Metrics are displayed to three decimals from locked archived JSON reports and include all SC-FMA variants available in the locked package. ``First'' denotes whether the highest-priority step is a first-selected target. N/C for the Simple average row indicates that first-selected target and Mass@25\% were not computed for this negatively correlated reference view.}
\label{tab:prm800k-audit}
\centering
\scriptsize
\renewcommand{\arraystretch}{1.08}
\begin{tabular}{Z{0.21\linewidth}ccccZ{0.20\linewidth}}
\toprule
\textbf{Field} & \textbf{PRM $\rho$} & \textbf{First} & \textbf{Mass@25\%} & \textbf{NDCG@25\%} & \textbf{Role} \\
\midrule
\wstruct{} & \textbf{0.611} & \textbf{0.917} & \textbf{0.381} & \textbf{0.951} & Strongest real-data fidelity field. \\
SC-FMA Ridge & 0.604 & 0.905 & 0.380 & 0.946 & Close fidelity-tracking decomposition; no $\gamma/\delta$ terms. \\
SC-FMA QP & 0.442 & 0.903 & 0.379 & 0.944 & Full redundancy/bottleneck view; not the PRM800K winner. \\
SC-FMA Projection & $-$0.135 & 0.723 & 0.288 & 0.743 & Lower-consistency projection view. \\
Frozen PRM prefix signal & 0.252 & 0.777 & 0.342 & 0.847 & In-distribution reference context. \\
Raw local utility & $-$0.077 & 0.686 & 0.282 & 0.722 & Direct ablation. \\
Simple average (local utility + necessity) & $-$0.443 & N/C & N/C & 0.438 & Independent feature-average reference. \\
Random & 0.006 & 0.733 & 0.307 & 0.776 & Seeded sanity control. \\
\bottomrule
\end{tabular}
\end{table}

Table~\ref{tab:prm800k-audit} shows annotation-order behavior on a process-annotation distribution. The normalized fidelity field $\tilde{\mathbf{c}}$ is instantiated by the supervised \wstruct{} field, and Ridge closely tracks this field while adding decomposition fields. The archived fixed-budget readouts show a small tracking cost (Mass@25\% 0.3809 for \wstruct{} versus 0.3796 for Ridge; Coverage NDCG@25\% 0.9506 versus 0.9460). The simple-average reference view (unweighted mean of raw local utility and structural necessity, Spearman $\rho = -0.443$) loses annotation-order consistency. Its component audit explains the sign: raw local utility is weakly negative ($\rho=-0.077$), and the structural-necessity field alone is more strongly anti-correlated with the PRM800K labels ($\rho=-0.418$). This negative direction means the current temporal-edge and TF-IDF graph construction is not merely low-yield in this domain; it is directionally misaligned with the human process labels under this readout. Thus, in this real-data route, naive averaging of raw utility with structural necessity is not trustworthy; a learned fidelity field and a conservative Ridge combination are needed when the fidelity field is already strong. Because no same-supervision structure-only Ridge control is reported, the comparison cannot establish that a supervised structural model would fail under equal supervision.

\subsection{Representation Enrichment Analysis}
\label{sec:representation-enrichment}

The close correlation between \wstruct{} and SC-FMA Ridge makes the enrichment analysis central. We compare three views already stored or derivable in the locked audit-card auto-validation output: a scalar \wstruct{} representation, a raw-field bundle that exposes fidelity, raw local utility, structural necessity, redundancy density, and bottleneck indicators without SCU joint optimization, and the structured SC-FMA view with recommended actions and utility-anchor interpretation fields.

\textbf{Field enrichment.} The scalar baseline exposes one field: \texttt{w\_struct\_scalar}. The raw-field bundle exposes five unoptimized fields: fidelity, raw local utility, structural necessity, redundancy density, and bottleneck indicator. The SC-FMA view exposes six fields: fidelity, structural necessity, redundancy, bottleneck status, recommended action, and upstream utility-anchor status. These fields operationalize dependency, redundancy, bottleneck, structural-role, audit-reason, and utility-anchor dimensions. The raw-field bundle is included to test whether the retrieval gain comes merely from exposing more fields.

\begin{table}[ht]
\caption{Representation enrichment relative to a scalar \wstruct{} field. Counts are taken from the locked audit-card auto-validation output and report additional audit-record fields and audit categories.}
\label{tab:representation-enrichment}
\centering
\scriptsize
\renewcommand{\arraystretch}{1.08}
\begin{tabular}{Z{0.24\linewidth}cccccc r}
\toprule
\textbf{Representation} & \textbf{Fid.} & \textbf{Dep.} & \textbf{Red.} & \textbf{Bot.} & \textbf{Reason} & \textbf{Anchor} & \textbf{Fields} \\
\midrule
\wstruct{} scalar & \checkmark & -- & -- & -- & -- & -- & 1 \\
Raw field bundle & \checkmark & \checkmark & \checkmark & \checkmark & -- & \checkmark & 5 \\
SC-FMA record & \checkmark & \checkmark & \checkmark & \checkmark & \checkmark & \checkmark & 6 \\
\bottomrule
\end{tabular}
\end{table}

The same output records four unique target categories: weak utility anchor, structural over-correction, redundancy consolidation, and bottleneck protection. Under the fixed 25\% inspection budget, the raw-field bundle raises mean rule-target recall from 0.235 to 0.524 and mean NDCG from 0.353 to 0.768, showing that field exposure explains a large share of the automatic retrieval signal. The SC-FMA view further raises mean recall to 0.699 and mean NDCG to 0.978. The target-dependence caveats are defined in Section~\ref{sec:why-calibration}; here, the readout is used only as audit-category separation, and the main lowest-overlap subset is analyzed separately below.

\begin{table}[ht]
\caption{Audit-category separation under the fixed 25\% audit budget. The targets are rule-derived categories in the locked audit-card auto-validation output; see Section~\ref{sec:why-calibration} for target-dependence caveats.}
\label{tab:audit-category-separation}
\centering
\small
\renewcommand{\arraystretch}{1.08}
\begin{tabular}{Z{0.27\linewidth}rrrr}
\toprule
\textbf{Audit category} & \textbf{Positive cases} & \textbf{\wstruct{} recall} & \textbf{Raw-field recall} & \textbf{SC-FMA recall} \\
\midrule
Weak utility anchor & 2,420 & 0.247 & 0.455 & 0.455 \\
Structural over-correction & 2,145 & 0.228 & 0.430 & 0.515 \\
Redundancy consolidation & 1,214 & 0.199 & 0.762 & 0.824 \\
Bottleneck protection & 382 & 0.264 & 0.450 & 1.000 \\
\bottomrule
\end{tabular}
\end{table}

\textbf{Representation disagreement cases.} At the locked-summary level, Ridge is close to \wstruct{} on annotation-order consistency (0.604 versus 0.611), while QP is more distinct (0.442). Ridge is intentionally fidelity-tracking-oriented. QP shows how the summary order changes when the record shifts from fidelity tracking toward constrained structural allocation. The archived package stores summary readouts for these variants; per-artifact top-$k$ overlap and rank-change ratios are outside the stored output schema.

\textbf{Audit-category separation.} The clearest enrichment appears when artifacts or traces have similar scalar-fidelity visibility but different audit reasons. Under the same 25\% budget, raw fields already separate many weak utility-anchor and redundancy cases that the scalar view collapses. SC-FMA adds beyond raw fields where a record-level interpretation is needed, especially structural over-correction and redundancy consolidation. The target-dependence caveats are handled in Section~\ref{sec:why-calibration}. These readouts support the representation-level claim that SC-FMA increases audit organization, not the stronger claim that SCU optimization alone explains all retrieval gains.

\subsection{Knowledge-Graph Backend Feasibility Study}
\label{sec:kg-pilot}

The Countries-KG backend feasibility study examines typed graph ingestion within the audit-record interface. As a check on non-lexical graph construction, we use the Countries KG resource associated with approximate-reasoning evaluation over country, region, and neighborhood relations~\cite{bouchard2015approximate}. The fixture contains 30 entities and 189 triples over \texttt{locatedIn} and \texttt{neighborOf}. We generate 30 deterministic KG-query traces and construct both a lightweight lexical-dependency graph and a KG-augmented graph that maps ontology relations to functional edge types.

\begin{table}[ht]
\caption{Countries-KG backend feasibility study (30 traces over a real knowledge graph, 30 entities, 189 triples). KG-augmented edges are derived from typed ontology relations. Structural-necessity consistency is the archived rank-consistency check between TF-IDF reference and KG-augmented necessity fields; the reported 0.0 is the rounded archived value under that computation.}
\label{tab:kg-pilot}
\centering
\small
\renewcommand{\arraystretch}{1.08}
\begin{tabular}{Z{0.45\linewidth}Z{0.43\linewidth}}
\toprule
\textbf{Quantity} & \textbf{Value} \\
\midrule
KG entities / triples & 30 / 189 \\
KG relation types & \texttt{locatedIn}, \texttt{neighborOf} \\
Reasoning traces & 30 \\
TF-IDF reference edges & 310 \\
KG candidate edges & 304 \\
KG-augmented edges & 320 \\
KG-derived edges added & 196 \\
Ontology relations detected & depends\_on\_concept (211), same\_constraint (85), neighbor\_concept (8) \\
Mean abs structural-necessity delta & 0.431 \\
Structural-necessity consistency (reference vs KG) & 0.0 \\
Bottleneck top-1 overlap & 1.000 \\
Matched nodes & 154 \\
\bottomrule
\end{tabular}
\end{table}

Table~\ref{tab:kg-pilot} shows that typed KG edges change the graph substantially: the stage considers 304 KG candidate relations, adds 196 ontology-derived edges, and yields 320 KG-augmented graph edges versus 310 TF-IDF reference edges. The mean absolute structural-necessity delta is 0.431 with an archived reference-vs-KG necessity consistency rounded to 0.0. One reading is that KG augmentation changes which objects the representation marks as structurally dependent. The equally important caution is that the structural-necessity field is highly sensitive to the graph-construction backend in this fixture; this sensitivity is a reliability warning, not a validation win. The 30-trace fixture is proof-of-concept evidence for typed backend substitution and a methodological analogy to knowledge-graph curation: the evidence concerns the graph-construction interface, not production-scale ontology validation or a live knowledge-base audit workflow. Full edge lists, the audit report, and degeneracy-handling details are in the supplementary material.
% Scope note for tests: fixture-level typed-edge construction only; no deployed workflow validation.

\subsection{Rule-Derived Audit Record Retrieval Check}
\label{sec:why-calibration}
\label{sec:audit-results}
\label{sec:oracle-auto-validation}

The audit-target retrieval stage examines whether decomposition fields make rule-derived audit targets visible under a fixed review budget. The main readout is the weak utility-anchor family, which has the lowest direct overlap with the exposed SC-FMA output fields. In this stage, \wstruct{} is the frozen scalar fidelity field supplied to the record builder. The weak utility-anchor rule uses this input-side fidelity-field repair pattern; redundancy and structural over-correction targets reuse fields that SC-FMA exposes, and bottleneck protection directly reuses QP budget entry. The raw-field bundle is therefore reported alongside SC-FMA to separate field exposure from SCU joint optimization.

We evaluate failure-mode localization as retrieval over traces. The scalar-only view receives the \wstruct{} summary, the raw-field view receives unoptimized primitive fields, and the SC-FMA view receives fidelity, necessity, redundancy, bottleneck status, recommended-action, and interpretation fields. Automatic oracle targets come from frozen process-annotation labels, priority-field disagreements, redundancy density, bottleneck flags, and the failure-taxonomy rules in Appendix~\ref{app:failure-taxonomy}. Weak utility-anchor targets have the lowest direct field overlap; redundancy consolidation and structural over-correction provide partially overlapping audit context. Bottleneck protection is QP-derived by construction and is treated separately as a record-consistency check.

The target-family rules are fixed before scoring and are implemented as explicit threshold tests. In compact form: structural over-correction is assigned when $\rho_{\mathrm{QP}}\le \rho_{\wstruct{}} - 0.10$; redundancy consolidation when redundancy density is in the top tertile and QP trails $\max(\rho_{\mathrm{Ridge}},\rho_{\wstruct{}})$ by at least 0.05; bottleneck protection when a bottleneck step enters the QP top-25\% budget while its process label is below the trace median; and weak utility anchor when $\rho_{\mathrm{raw}}\le 0$ and $\rho_{\wstruct{}}\ge 0.30$. These labels are multi-label and not mutually exclusive. Because the bottleneck-protection definition directly uses QP budget entry, its SC-FMA score checks whether the audit record reflects an internally computed QP bottleneck state. It is an internal record-consistency check rather than independent target-retrieval validation. Weak utility-anchor is the lowest direct-overlap rule, but it is still not fully independent because it uses \wstruct{} as an input-side fidelity field.

\begin{table}[ht]
\caption{Field-overlap audit for rule-derived target families. Overlap is the number of SC-FMA decomposition fields directly reused by the target rule divided by the six fields exposed in the archived SC-FMA view.}
\label{tab:target-circularity-audit}
\centering
\scriptsize
\renewcommand{\arraystretch}{1.08}
\begin{tabular}{Z{0.21\linewidth}Z{0.25\linewidth}cZ{0.22\linewidth}Z{0.14\linewidth}}
\toprule
\textbf{Target family} & \textbf{Fields used by target rule} & \textbf{Overlap} & \textbf{SC-FMA field relation} & \textbf{Interpretation} \\
\midrule
Weak utility anchor & Process labels, raw local utility, frozen \wstruct{} fidelity-field repair pattern & 1/6 & Lowest direct overlap; uses an input-side fidelity repair pattern & Main low-overlap retrieval check \\
Redundancy consolidation & Redundancy density, redundant-cluster flags, priority change & 2/6 & Partial overlap through redundancy and recommended-action fields & Exploratory check \\
Bottleneck protection & Bottleneck flag, QP top-25\% budget entry, below-median process label & 2/6 & QP-derived through bottleneck and budget-entry fields & Internal consistency check \\
Structural over-correction & Fidelity-vs-structure disagreement, QP/Ridge priority divergence & 2/6 & Partial overlap through fidelity and necessity disagreement fields & Exploratory check \\
\bottomrule
\end{tabular}
\end{table}

\begin{table}[ht]
\caption{Per-target audit-target retrieval point estimates under the same 25\% budget. Field overlap reports the approximate definition overlap with the six SC-FMA decomposition fields. Asterisks mark target families with stronger definitional overlap; see Section~\ref{sec:why-calibration} for the weak utility-anchor and bottleneck-protection caveats.}
\label{tab:oracle-auto-validation-by-target}
\centering
\scriptsize
\renewcommand{\arraystretch}{1.08}
\begin{tabular}{Z{0.21\linewidth}rZ{0.09\linewidth}cccccc}
\toprule
\textbf{Target family} & \textbf{Pos.} & \textbf{Overlap} & \textbf{\wstruct{} R} & \textbf{Raw R} & \textbf{SC-FMA R} & \textbf{\wstruct{} N} & \textbf{Raw N} & \textbf{SC-FMA N} \\
\midrule
Bottleneck protection\textsuperscript{*} & 382 & 2/6 QP & 0.264 & 0.450 & 1.000 & 0.221 & 0.411 & 1.000 \\
Redundancy consolidation\textsuperscript{*} & 1,214 & 2/6 & 0.199 & 0.762 & 0.824 & 0.224 & 0.849 & 0.916 \\
Weak utility anchor & 2,420 & 1/6 input & 0.247 & 0.455 & 0.455 & 0.527 & 0.997 & 0.997 \\
Structural over-correction\textsuperscript{*} & 2,145 & 2/6 & 0.228 & 0.430 & 0.515 & 0.439 & 0.817 & 1.000 \\
\bottomrule
\end{tabular}
\end{table}

The weak utility-anchor row provides the lowest-direct-overlap subset check in the locked output and is the main audit-target retrieval readout. The raw-field and SC-FMA rows are identical on this target, so this subset does not show an SCU-specific retrieval advantage. The rule uses process labels and the frozen \wstruct{} fidelity field from the same archived package. The input-output distinction is important: \wstruct{} supplies the scalar fidelity field for this stage, raw fields expose primitive components, and the SC-FMA output adds record-level dependency, redundancy, bottleneck, reason, and interpretation fields. The per-target rows are point estimates; the archived aggregate report stores bootstrap intervals only for the aggregate retrieval table below. Table~\ref{tab:oracle-auto-validation} reports aggregate automatic retrieval as contextual evidence because targets with definitional overlap account for 3,741 of 6,161 target positions (about 61\%).

\begin{table}[ht]
\caption{Aggregate retrieval under partially overlapping target definitions. Failure-mode localization is evaluated on the locked PRM800K split as a representation-coverage check under a 25\% budget. The scalar condition uses \wstruct{}, the raw-field bundle exposes primitive fields without SCU joint optimization, and the SC-FMA condition uses decomposition and interpretation fields. ``First'' denotes whether the first returned target appears under the fixed budget; target-dependence caveats are given above and in Table~\ref{tab:oracle-auto-validation-by-target}.}
\label{tab:oracle-auto-validation}
\centering
\scriptsize
\renewcommand{\arraystretch}{1.08}
\begin{tabular}{Z{0.30\linewidth}ccccc}
\toprule
\textbf{View} & \textbf{Cov.@25\%} & \textbf{NDCG@25\%} & \textbf{MRR} & \textbf{First} & \textbf{Cost} \\
\midrule
\wstruct{} scalar only & 0.235 & 0.353 & 0.354 & 0.000 & 0.0007 \\
Raw field bundle & 0.524 & 0.768 & 1.000 & 1.000 & 0.0002 \\
SC-FMA decomposition & 0.699 & 0.978 & 1.000 & 1.000 & 0.0002 \\
\bottomrule
\end{tabular}
\end{table}

The aggregate difference is exploratory recovery of rule-defined audit targets by a decomposed audit record. The archived aggregate bootstrap intervals are wide for coverage and NDCG: \wstruct{} coverage $[0.211,0.256]$ and NDCG $[0.222,0.483]$, raw-field coverage $[0.435,0.684]$ and NDCG $[0.521,0.952]$, and SC-FMA coverage $[0.485,0.912]$ and NDCG $[0.937,1.000]$. These intervals do not remove the definitional-overlap concern. The check tests whether the record fields are machine-readable; independent audit quality remains a human-validation question. The raw-field control shows that much of the gain comes from exposing primitive fields, while SC-FMA adds record-level organization and additional retrieval mainly for redundancy and structural over-correction families. The bottleneck contribution follows the internal-consistency interpretation above. Table~\ref{tab:compact-audit-card-case} gives one compact preview case; the appendix Audit Cards provide fuller trace-level examples.

\begin{table}[H]
\caption{Compact Audit Card case study: scalar-only view versus SC-FMA decomposition.}
\label{tab:compact-audit-card-case}
\centering
\small
\renewcommand{\arraystretch}{1.08}
\begin{tabular}{Z{0.22\linewidth}Z{0.30\linewidth}Z{0.36\linewidth}}
\toprule
\textbf{Failure target} & \textbf{\texttt{wstruct}-only view} & \textbf{SC-FMA Audit Card view} \\
\midrule
Bottleneck protection & A scalar signal can allocate a step for early inspection while leaving downstream gating implicit. & Bottleneck and necessity fields mark the step as dependency-exposed, so the interpretation is to inspect the downstream chain before dropping it. \\
Redundancy consolidation & Multiple high-signal steps can look independently important. & Redundancy density and neighborhood fields identify a cluster, so the interpretation is to consolidate the checks or sample the cluster. \\
Weak utility anchor & A high calibrated fidelity field can hide a failed raw local utility signal. & Fidelity and raw-anchor disagreement expose an upstream utility-anchor problem, so the interpretation is to diagnose the signal source. \\
\bottomrule
\end{tabular}
\end{table}

\subsection{Controlled Calibration Analysis}
\label{sec:controlled-calibration}

\begin{table}[ht]
\caption{Controlled representation-field assignment on the synthetic stage. For SC-FMA variants, Spearman $\rho$ is the 5-seed mean $\pm$ standard deviation; remaining columns are from representative seed 42 because the archived multiseed run records Spearman alone. Simple average is the element-wise mean of normalized local utility and structural necessity.}
\label{tab:controlled-calibration-results}
\centering
\small
\begin{tabular}{Z{0.30\linewidth}cccc}
\toprule
\textbf{View} & \textbf{Synthetic label $\rho$} & \textbf{Kendall $\tau$} & \textbf{Record NDCG@3} & \textbf{Record NDCG@5} \\
\midrule
SC-FMA QP                  & \textbf{0.597} $\pm$ 0.064 & 0.533 & 0.921 & 0.958 \\
SC-FMA Ridge               & 0.541 $\pm$ 0.016 & 0.447 & 0.889 & 0.927 \\
Raw local utility          & 0.483 & 0.412 & 0.874 & 0.919 \\
Simple average (local utility + necessity) & 0.478 & 0.413 & 0.889 & 0.926 \\
SC-FMA Projection          & 0.372 $\pm$ 0.030 & 0.278 & 0.824 & 0.890 \\
Span Length                & $-$0.006 & $-$0.013 & 0.740 & 0.826 \\
Relative Position          & $-$0.005 & $-$0.004 & 0.735 & 0.827 \\
Random                     & $-$0.010 & $-$0.008 & 0.736 & 0.825 \\
\bottomrule
\end{tabular}
\end{table}

Table~\ref{tab:controlled-calibration-results} shows the intended calibration behavior: SC-FMA QP records higher annotation-order consistency than raw local utility, with a Spearman gap of 0.114. Under the controlled annotation target, structural calibration improves local-utility field assignment.

\begin{table}[ht]
\caption{SCU component contribution on the controlled synthetic stage. Values are 5-seed mean $\pm$ standard deviation over seeds 42, 123, 456, 789, and 1024. Delta is the annotation-order consistency difference relative to \texttt{full\_scu}. The largest drop occurs when the fidelity field is removed; removing redundancy or bottleneck regularization has near-zero impact on the reported consistency readouts in this setting.}
\label{tab:scu-component-contribution}
\centering
\small
\renewcommand{\arraystretch}{1.08}
\begin{tabular}{lccc}
\toprule
\textbf{Variant} & \textbf{Synthetic label $\rho$ mean$\pm$std} & \textbf{Kendall $\tau$ mean$\pm$std} & \textbf{Delta vs Full} \\
\midrule
\texttt{full\_scu} & 0.5103 $\pm$ 0.0320 & 0.4297 $\pm$ 0.0287 & 0.0000 \\
\texttt{no\_fidelity} & 0.2850 $\pm$ 0.0153 & 0.2366 $\pm$ 0.0138 & $-$0.2253 \\
\texttt{no\_structure} & 0.4841 $\pm$ 0.0334 & 0.4053 $\pm$ 0.0285 & $-$0.0263 \\
\texttt{no\_redundancy} & 0.5093 $\pm$ 0.0321 & 0.4273 $\pm$ 0.0273 & $-$0.0010 \\
\texttt{no\_bottleneck} & 0.5090 $\pm$ 0.0304 & 0.4279 $\pm$ 0.0274 & $-$0.0013 \\
\bottomrule
\end{tabular}
\end{table}

Table~\ref{tab:scu-component-contribution} refines the component narrative. Removing fidelity reduces annotation consistency by 0.2253, far more than any other removal. Graph necessity contributes moderately (Delta $=-0.0263$), while redundancy and bottleneck change consistency by about 0.001. Here, those terms mainly provide regularization and decomposition. For annotation-order consistency alone, the simplified Ridge view is often sufficient; redundancy and bottleneck fields matter most when the audit task asks why a selected artifact is duplicated, dependency-exposed, or protected under a fixed budget.

The near-zero redundancy and bottleneck deltas reflect the structure-independent proxy label. We therefore use a structural stress-test (Table~\ref{tab:scu-stress-test}) whose labels encode redundancy-awareness and bottleneck protection.

\begin{table}[ht]
\caption{SCU component contribution on a structural stress-test fixture (200 traces per seed, five fixed seeds 42, 123, 456, 789, 1024). These designed labels reward redundancy-awareness and bottleneck protection under a stress-test redundancy setting ($\gamma=0.6$). Delta is the annotation-order consistency difference relative to \texttt{full\_scu}. The table is an implementation and mechanism check, not real-data audit-value evidence.}
\label{tab:scu-stress-test}
\centering
\small
\renewcommand{\arraystretch}{1.08}
\begin{tabular}{lccc}
\toprule
\textbf{Variant} & \textbf{Stress-label $\rho$ mean$\pm$std} & \textbf{Kendall $\tau$ mean$\pm$std} & \textbf{Delta vs Full} \\
\midrule
\texttt{full\_scu} & 0.6249 $\pm$ 0.0145 & 0.5238 $\pm$ 0.0112 & 0.0000 \\
\texttt{no\_fidelity} & 0.4832 $\pm$ 0.0357 & 0.3935 $\pm$ 0.0319 & $-$0.1417 \\
\texttt{no\_structure} & 0.5871 $\pm$ 0.0234 & 0.4870 $\pm$ 0.0224 & $-$0.0378 \\
\texttt{no\_redundancy} & 0.5758 $\pm$ 0.0163 & 0.4889 $\pm$ 0.0091 & $-$0.0492 \\
\texttt{no\_bottleneck} & 0.5409 $\pm$ 0.0145 & 0.4424 $\pm$ 0.0114 & $-$0.0840 \\
\bottomrule
\end{tabular}
\end{table}

On the stress-test, removing fidelity costs the most ($-0.1417$). Structural terms also contribute: bottleneck log-barrier $0.0840$, redundancy penalty $0.0492$, and structural-dependency term $0.0378$. Redundancy and bottleneck act as conditional regularizers whose contribution grows when the fidelity field lacks structural preference or when high-redundancy blocks make audit interpretation depend on cluster structure.

\begin{table}[ht]
\caption{Graph construction analysis on the development run. Temporal edges define the primary backbone; topical similarity edges are supplementary cues. Record NDCG@3 and Record NDCG@5 are interpreted as budget-constrained audit-coverage readouts computed from the same source data and metric definition used for the correlation readout.}
\label{tab:graph-construction-analysis}
\centering
\small
\renewcommand{\arraystretch}{1.08}
\begin{tabular}{lcccc}
\toprule
\textbf{Variant} & \textbf{Graph-check $\rho$} & \textbf{Kendall $\tau$} & \textbf{Record NDCG@3} & \textbf{Record NDCG@5} \\
\midrule
\texttt{full\_tfidf} & 0.0515 & 0.0458 & 0.8655 & 0.8655 \\
\texttt{temporal\_only} & 0.0596 & 0.0527 & 0.8688 & 0.8688 \\
\texttt{jaccard\_topical} & 0.0517 & 0.0460 & 0.8657 & 0.8657 \\
\texttt{shuffled\_topical} & 0.0515 & 0.0458 & 0.8655 & 0.8655 \\
\bottomrule
\end{tabular}
\end{table}

Table~\ref{tab:graph-construction-analysis} shows that topical edges add little annotation-signal tracking evidence: \texttt{full\_tfidf} is slightly below \texttt{temporal\_only} and matches \texttt{shuffled\_topical}. Their observed role is supplementary structural cueing.

The reference-view pattern clarifies the source of the representation signal. Raw local utility carries more annotation-order information than span length, position, and random controls but is still negative on the locked process-annotation route. Plain averaging of local utility and structural necessity is worse than either conservative \wstruct{} or Ridge because the structural-necessity input is not independently aligned with PRM800K labels in that setting. The archived workflow is defined at the artifact-step level.

The component and graph-construction analyses clarify how the representation is built. SCU benefits primarily from fidelity-field tracking and only conditionally from graph-necessity alignment. On the PRM800K route, graph necessity is directionally misaligned with human labels; on the designed stress-test, structural terms help because the labels encode the structural preference by construction. Redundancy and bottleneck terms mainly support conservative decomposition. Topical graph edges move structural readouts while leaving annotation-order consistency largely unchanged.

Error analysis identifies three residual patterns: high-utility but zero-necessity steps can receive excessive mass; near-threshold bottlenecks can be under-protected; and dense redundancy blocks can over-equalize priority values.

\subsection{Failure Modes of Representation Layer}
\label{sec:prm800k-error}

The stratified audit-record checks show when structural calibration helps or hinders process-annotation consistency. The multi-label failure taxonomy over 4,417 locked traces reports weak utility anchor in 54.79\% of traces, structural over-correction in 48.56\%, redundancy misclassification in 27.48\%, low signal or tie in 18.32\%, and bottleneck over-protection in 8.65\%. Fixed-template Audit Cards are provided in Appendix~\ref{app:failure-taxonomy}.

Trace length is an important observed predictor of variant behavior and a practical limitation for the full QP variant. On short traces (7,141 steps across 1,806 samples), \wstruct{} reaches Spearman $\rho = 0.759$, Ridge closely tracks it at 0.757, and QP remains close at 0.734. The middle trace-length stratum (8,264 steps across 1,195 samples) is intermediate: \wstruct{} is 0.583, Ridge is 0.579, and QP is 0.321. On long traces (18,814 steps, 1,416 samples), \wstruct{} drops to 0.447, Ridge to 0.430, and QP degrades to 0.172. Longer and more complex traces are exactly where audit support is often needed, so this degradation is operationally important rather than a minor tail case. Longer traces likely create denser redundancy blocks; QP's hard penalty then reallocates the audit-priority field away from the annotation signal, while Ridge's soft averaging avoids the sharpest over-correction. This pattern motivates a future engineering extension: sparse, blockwise, or windowed calibration could reduce the effective $k$ in the constrained subproblem while preserving the same SCU record fields.

Label entropy gives a complementary view. On low-entropy traces (17,288 steps, 1,474 samples), discriminability is limited but \wstruct{} and Ridge still record Coverage NDCG@25\% above 0.99 as a fixed-budget audit-coverage readout. The middle entropy stratum (10,948 steps across 1,639 samples) gives \wstruct{} 0.642, Ridge 0.632, and QP 0.501. On high-entropy traces (5,983 steps, 1,304 samples), labels are more dispersed and QP ($\rho = 0.642$) moves closer to \wstruct{} ($\rho = 0.705$), suggesting that explicit redundancy management is more useful when annotation structure is richer.

Across strata, \wstruct{} provides the main fidelity reference, Ridge stays within 0.01--0.02 Spearman $\rho$ across most strata, QP degrades on long or high-redundancy traces, and Projection provides the least consistent SC-FMA representation variant. The frozen PRM prefix signal provides moderate in-distribution reference context (Spearman 0.061--0.425) and remains consistently above raw local utility.

